\documentclass[12pt]{article}

% Idioma y codificación
\usepackage[spanish, es-tabla, es-notilde]{babel}       %es-tabla para que se titule "Tabla"
\usepackage[utf8]{inputenc}

% Márgenes
\usepackage[a4paper,top=3cm,bottom=2.5cm,left=3cm,right=3cm]{geometry}

% Comentarios de bloque
\usepackage{verbatim}

% Paquetes de links
\usepackage[hidelinks]{hyperref}    % Permite enlaces
\usepackage{url}                    % redirecciona a la web

% Más opciones para enumeraciones
\usepackage{enumitem}

% Personalizar la portada
\usepackage{titling}

% Paquetes de tablas
\usepackage{multirow}

% Para añadir el símbolo de euro
\usepackage{eurosym}

% Paquete para subrayar con diferentes estilos
\usepackage[normalem]{ulem}


%------------------------------------------------------------------------

%Paquetes de figuras
\usepackage{caption}
\usepackage{subcaption} % Figuras al lado de otras
\usepackage{float}      % Poner figuras en el sitio indicado H.


% Paquetes de imágenes
\usepackage{graphicx}       % Paquete para añadir imágenes
\usepackage{transparent}    % Para manejar la opacidad de las figuras

% Paquete para usar colores
\usepackage[dvipsnames, table, xcdraw]{xcolor}
\usepackage{pagecolor}      % Para cambiar el color de la página

% Habilita tamaños de fuente mayores
\usepackage{fix-cm}

% Para los gráficos
\usepackage{tikz}
\usepackage{forest}

% Para poder situar los nodos en los grafos
\usetikzlibrary{positioning}

% Para evitar problemas entre babel y tikz
\usetikzlibrary{babel}


%------------------------------------------------------------------------

% Paquetes de matemáticas
\usepackage{mathtools, amsfonts, amssymb, mathrsfs}
\usepackage[makeroom]{cancel}     % Simplificar tachando
\usepackage{polynom}    % Divisiones y Ruffini
\usepackage{units} % Para poner fracciones diagonales con \nicefrac

\usepackage{pgfplots}   %Representar funciones
\pgfplotsset{compat=1.18}  % Versión 1.18

\usepackage{tikz-cd}    % Para usar diagramas de composiciones
\usetikzlibrary{calc}   % Para usar cálculo de coordenadas en tikz

%Definición de teoremas, etc.
\usepackage{amsthm}
%\swapnumbers   % Intercambia la posición del texto y de la numeración

\theoremstyle{plain}

\makeatletter
\@ifclassloaded{article}{
  % --- Configuración para ARTICLE ---
  \newtheorem{teo}{Teorema}[section]
  \theoremstyle{definition}
  \newtheorem{definicion}{Definición}[section]
  \newtheorem{ejercicio}{Ejercicio}[section]

  \renewcommand{\theteo}{%
    \ifnum\value{section}=0 \arabic{teo}\else \thesection.\arabic{teo}\fi
  }
  \renewcommand{\thedefinicion}{%
    \ifnum\value{section}=0 \arabic{definicion}\else \thesection.\arabic{definicion}\fi
  }
  \renewcommand{\theejercicio}{%
    \ifnum\value{section}=0 \arabic{ejercicio}\else \thesection.\arabic{ejercicio}\fi
  }
}{
  % --- Configuración para BOOK / REPORT ---
  \newtheorem{teo}{Teorema}[chapter]
  \theoremstyle{definition}
  \newtheorem{definicion}{Definición}[chapter]
  \newtheorem{ejercicio}{Ejercicio}[section]

  % Si no hay sección: Capítulo.Teorema (ej. 2.1)
  % Si hay sección:    Capítulo.Sección.Teorema (ej. 2.3.1)
  \renewcommand{\theteo}{%
    \ifnum\value{section}=0
      \thechapter.\arabic{teo}%
    \else
      \thesection.\arabic{teo}% (\thesection ya contiene el número de capítulo)
    \fi
  }
  \renewcommand{\thedefinicion}{%
    \ifnum\value{section}=0
      \thechapter.\arabic{definicion}%
    \else
      \thesection.\arabic{definicion}%
    \fi
  }
  \renewcommand{\theejercicio}{%
    \ifnum\value{section}=0
      \arabic{ejercicio}%
    \else
      \thesection.\arabic{ejercicio}% Formato por defecto: ej. 2.2.1
    \fi
  }
}
\makeatother

% Teoremas dependientes o con contador compartido
\newtheorem{coro}{Corolario}[teo]
\newtheorem{prop}[teo]{Proposición}
\newtheorem{lema}[teo]{Lema}
\newtheorem*{lema*}{Lema}

\theoremstyle{remark}
\newtheorem*{observacion}{Observación}

\theoremstyle{definition}
\newtheorem*{notacion}{Notación}
\newtheorem*{ejemplo}{Ejemplo}
\newtheorem*{ejercicio*}{Ejercicio}


% \renewcommand\qedsymbol{$\blacksquare$}         % Cambiar símbolo QED
%------------------------------------------------------------------------

% Paquetes para encabezados
\usepackage{fancyhdr}
\pagestyle{fancy}
\fancyhf{}

\newcommand{\helv}{ % Modificación tamaño de letra
\fontfamily{}\fontsize{12}{12}\selectfont}
\setlength{\headheight}{15pt} % Amplía el tamaño del índice


%\usepackage{lastpage}   % Referenciar última pag   \pageref{LastPage}
%\fancyfoot[C]{%
%  \begin{minipage}{\textwidth}
%    \centering
%    ~\\
%    \thepage\\
%    \href{https://losdeldgiim.github.io/}{\texttt{\footnotesize losdeldgiim.github.io}}
%  \end{minipage}
%}
\fancyfoot[C]{\thepage}
\fancyfoot[R]{\href{https://losdeldgiim.github.io/}{\texttt{\footnotesize losdeldgiim.github.io}}}

%------------------------------------------------------------------------

% Conseguir que no ponga "Capítulo 1". Sino solo "1."
\makeatletter
\@ifclassloaded{book}{
  \renewcommand{\chaptermark}[1]{\markboth{\thechapter.\ #1}{}} % En el encabezado
    
  \renewcommand{\@makechapterhead}[1]{%
  \vspace*{50\p@}%
  {\parindent \z@ \raggedright \normalfont
    \ifnum \c@secnumdepth >\m@ne
      \huge\bfseries \thechapter.\hspace{1em}\ignorespaces
    \fi
    \interlinepenalty\@M
    \Huge \bfseries #1\par\nobreak
    \vskip 40\p@
  }}
}
\makeatother

%------------------------------------------------------------------------
% Paquetes de cógido
\usepackage{minted}
\renewcommand\listingscaption{Código fuente}

\usepackage{fancyvrb}
% Personaliza el tamaño de los números de línea
\renewcommand{\theFancyVerbLine}{\small\arabic{FancyVerbLine}}

% Estilo para C++
\newminted{cpp}{
    frame=lines,
    framesep=2mm,
    baselinestretch=1.2,
    linenos,
    escapeinside=||
}

% para minted
\definecolor{LightGray}{rgb}{0.95,0.95,0.92}
\setminted{
    linenos=true,
    stepnumber=5,
    numberfirstline=true,
    autogobble,
    breaklines=true,
    breakautoindent=true,
    breaksymbolleft=,
    breaksymbolright=,
    breaksymbolindentleft=0pt,
    breaksymbolindentright=0pt,
    breaksymbolsepleft=0pt,
    breaksymbolsepright=0pt,
    fontsize=\footnotesize,
    bgcolor=LightGray,
    numbersep=10pt
}


\usepackage{listings} % Para incluir código desde un archivo

\renewcommand\lstlistingname{Código Fuente}
\renewcommand\lstlistlistingname{Índice de Códigos Fuente}

% Definir colores
\definecolor{vscodepurple}{rgb}{0.5,0,0.5}
\definecolor{vscodeblue}{rgb}{0,0,0.8}
\definecolor{vscodegreen}{rgb}{0,0.5,0}
\definecolor{vscodegray}{rgb}{0.5,0.5,0.5}
\definecolor{vscodebackground}{rgb}{0.97,0.97,0.97}
\definecolor{vscodelightgray}{rgb}{0.9,0.9,0.9}

% Configuración para el estilo de C similar a VSCode
\lstdefinestyle{vscode_C}{
  backgroundcolor=\color{vscodebackground},
  commentstyle=\color{vscodegreen},
  keywordstyle=\color{vscodeblue},
  numberstyle=\tiny\color{vscodegray},
  stringstyle=\color{vscodepurple},
  basicstyle=\scriptsize\ttfamily,
  breakatwhitespace=false,
  breaklines=true,
  captionpos=b,
  keepspaces=true,
  numbers=left,
  numbersep=5pt,
  showspaces=false,
  showstringspaces=false,
  showtabs=false,
  tabsize=2,
  frame=tb,
  framerule=0pt,
  aboveskip=10pt,
  belowskip=10pt,
  xleftmargin=10pt,
  xrightmargin=10pt,
  framexleftmargin=10pt,
  framexrightmargin=10pt,
  framesep=0pt,
  rulecolor=\color{vscodelightgray},
  backgroundcolor=\color{vscodebackground},
}

%------------------------------------------------------------------------
% La etiqueda IA sirve para indicar a partir de donde le pasamos el preambulo a la IA para el escaneo
% IA

% Comandos definidos
\newcommand{\bb}[1]{\mathbb{#1}}
\newcommand{\cc}[1]{\mathcal{#1}}

% I prefer the slanted \leq
\let\oldleq\leq % save them in case they're every wanted
\let\oldgeq\geq
\renewcommand{\leq}{\leqslant}
\renewcommand{\geq}{\geqslant}

% Si y solo si
\newcommand{\sii}{\iff}

% MCD y MCM
\DeclareMathOperator{\mcd}{mcd}
\DeclareMathOperator{\mcm}{mcm}

% Signo
\DeclareMathOperator{\sgn}{sgn}

% Imagen
\DeclareMathOperator{\Ima}{Im}

% Aplicación identidad
\newcommand{\id}{\mathrm{Id}}

% Letras griegas
\newcommand{\eps}{\epsilon}
\newcommand{\veps}{\varepsilon}
\newcommand{\lm}{\lambda}

\newcommand{\ol}{\overline}
\newcommand{\ul}{\underline}
\newcommand{\wt}{\widetilde}
\newcommand{\wh}{\widehat}

\let\oldvec\vec
\renewcommand{\vec}{\overrightarrow}

% Derivadas parciales
\newcommand{\del}[2]{\frac{\partial #1}{\partial #2}}
\newcommand{\Del}[3]{\frac{\partial^{#1} #2}{\partial #3^{#1}}}
\newcommand{\deld}[2]{\dfrac{\partial #1}{\partial #2}}
\newcommand{\Deld}[3]{\dfrac{\partial^{#1} #2}{\partial #3^{#1}}}


\newcommand{\AstIg}{\stackrel{(\ast)}{=}}
\newcommand{\Hop}{\stackrel{L'H\hat{o}pital}{=}}

\newcommand{\red}[1]{{\color{red}#1}} % Para integrales, destacar los cambios.

% Método de integración
\newcommand{\MetInt}[2]{
    \left[\begin{array}{c}
        #1 \\ #2
    \end{array}\right]
}

% Declarar aplicaciones
% 1. Nombre aplicación
% 2. Dominio
% 3. Codominio
% 4. Variable
% 5. Imagen de la variable
\newcommand{\Func}[5]{
    \begin{equation*}
        \begin{array}{rrll}
            \displaystyle #1:& \displaystyle  #2 & \longrightarrow & \displaystyle  #3\\
               & \displaystyle  #4 & \longmapsto & \displaystyle  #5
        \end{array}
    \end{equation*}
}

%------------------------------------------------------------------------


\newcommand{\myhref}[2]{\href{#1}{\color{blue}\uline{#2}}}

\usetikzlibrary{decorations.markings}
\begin{document}

    % 1. Foto de fondo
    % 2. Título
    % 3. Encabezado Izquierdo
    % 4. Color de fondo
    % 5. Coord x del titulo
    % 6. Coord y del titulo
    % 7. Fecha

    
    % 1. Foto de fondo
% 2. Título
% 3. Encabezado Izquierdo
% 4. Color de fondo
% 5. Coord x del titulo
% 6. Coord y del titulo
% 7. Fecha
% 8. Autor

\newcommand{\portada}[8]{

    \portadaBase{#1}{#2}{#3}{#4}{#5}{#6}{#7}{#8}
    \portadaBook{#1}{#2}{#3}{#4}{#5}{#6}{#7}{#8}
}

\newcommand{\portadaExamen}[8]{

    \portadaBase{#1}{#2}{#3}{#4}{#5}{#6}{#7}{#8}
    \portadaArticle{#1}{#2}{#3}{#4}{#5}{#6}{#7}{#8}
}




\newcommand{\portadaBase}[8]{

    % Tiene la portada principal y la licencia Creative Commons
    
    % 1. Foto de fondo
    % 2. Título
    % 3. Encabezado Izquierdo
    % 4. Color de fondo
    % 5. Coord x del titulo
    % 6. Coord y del titulo
    % 7. Fecha
    % 8. Autor
    
    
    \thispagestyle{empty}               % Sin encabezado ni pie de página
    \newgeometry{margin=0cm}        % Márgenes nulos para la primera página
    
    
    % Encabezado
    \fancyhead[L]{\helv #3}
    \fancyhead[R]{\helv \nouppercase{\leftmark}}
    
    
    \pagecolor{#4}        % Color de fondo para la portada
    
    \begin{figure}[p]
        \centering
        \transparent{0.3}           % Opacidad del 30% para la imagen
        
        \includegraphics[width=\paperwidth, keepaspectratio]{../../_assets/#1}
    
        \begin{tikzpicture}[remember picture, overlay]
            \node[anchor=north west, text=white, opacity=1, font=\fontsize{60}{90}\selectfont\bfseries\sffamily, align=left] at (#5, #6) {#2};
            
            \node[anchor=south east, text=white, opacity=1, font=\fontsize{12}{18}\selectfont\sffamily, align=right] at (9.7, 3) {\href{https://losdeldgiim.github.io/}{\textbf{Los Del DGIIM}, \texttt{\footnotesize losdeldgiim.github.io}}};
            
            \node[anchor=south east, text=white, opacity=1, font=\fontsize{12}{15}\selectfont\sffamily, align=right] at (9.7, 1.8) {Doble Grado en Ingeniería Informática y Matemáticas\\Universidad de Granada};
        \end{tikzpicture}
    \end{figure}
    
    
    \restoregeometry        % Restaurar márgenes normales para las páginas subsiguientes
    \pagecolor{white}       % Restaurar el color de página
    
    
    \newpage
    \thispagestyle{empty}               % Sin encabezado ni pie de página
    \begin{tikzpicture}[remember picture, overlay]
        \node[anchor=south west, inner sep=3cm] at (current page.south west) {
            \begin{minipage}{0.5\paperwidth}
                \href{https://creativecommons.org/licenses/by-nc-nd/4.0/}{
                    \includegraphics[height=2cm]{../../_assets/Licencia.png}
                }\vspace{1cm}\\
                Esta obra está bajo una
                \href{https://creativecommons.org/licenses/by-nc-nd/4.0/}{
                    Licencia Creative Commons Atribución-NoComercial-SinDerivadas 4.0 Internacional (CC BY-NC-ND 4.0).
                }\\
    
                Eres libre de compartir y redistribuir el contenido de esta obra en cualquier medio o formato, siempre y cuando des el crédito adecuado a los autores originales y no persigas fines comerciales. 
            \end{minipage}
        };
    \end{tikzpicture}
    
    
    
    % 1. Foto de fondo
    % 2. Título
    % 3. Encabezado Izquierdo
    % 4. Color de fondo
    % 5. Coord x del titulo
    % 6. Coord y del titulo
    % 7. Fecha
    % 8. Autor


}


\newcommand{\portadaBook}[8]{

    % 1. Foto de fondo
    % 2. Título
    % 3. Encabezado Izquierdo
    % 4. Color de fondo
    % 5. Coord x del titulo
    % 6. Coord y del titulo
    % 7. Fecha
    % 8. Autor

    % Personaliza el formato del título
    \pretitle{\begin{center}\bfseries\fontsize{42}{56}\selectfont}
    \posttitle{\par\end{center}\vspace{2em}}
    
    % Personaliza el formato del autor
    \preauthor{\begin{center}\Large}
    \postauthor{\par\end{center}\vfill}
    
    % Personaliza el formato de la fecha
    \predate{\begin{center}\huge}
    \postdate{\par\end{center}\vspace{2em}}
    
    \title{#2}
    \author{\href{https://losdeldgiim.github.io/}{Los Del DGIIM, \texttt{\large losdeldgiim.github.io}}
    \\ \vspace{0.5cm}#8}
    \date{Granada, #7}
    \maketitle
    
    \tableofcontents
}




\newcommand{\portadaArticle}[8]{

    % 1. Foto de fondo
    % 2. Título
    % 3. Encabezado Izquierdo
    % 4. Color de fondo
    % 5. Coord x del titulo
    % 6. Coord y del titulo
    % 7. Fecha
    % 8. Autor

    % Personaliza el formato del título
    \pretitle{\begin{center}\bfseries\fontsize{42}{56}\selectfont}
    \posttitle{\par\end{center}\vspace{2em}}
    
    % Personaliza el formato del autor
    \preauthor{\begin{center}\Large}
    \postauthor{\par\end{center}\vspace{3em}}
    
    % Personaliza el formato de la fecha
    \predate{\begin{center}\huge}
    \postdate{\par\end{center}\vspace{5em}}
    
    \title{#2}
    \author{\href{https://losdeldgiim.github.io/}{Los Del DGIIM, \texttt{\large losdeldgiim.github.io}}
    \\ \vspace{0.5cm}#8}
    \date{Granada, #7}
    \thispagestyle{empty}               % Sin encabezado ni pie de página
    \maketitle
    \vfill
}
    \portadaExamen{ffccA4.jpg}{Topología II\\Examen VIII}{Topología II. Examen VIII}{MidnightBlue}{-8}{28}{2025}{José Juan Urrutia Milán}

    \begin{description}
        \item[Asignatura] Topología II.
        \item[Curso Académico] 2020/21.
        \item[Grado] Grado en Matemáticas.
        % \item[Grupo] Grupo Único.
        % \item[Profesor] José Antonio Gálvez.
        \item[Descripción] Convocatoria Extraordinaria.
        % \item[Fecha] 28 de noviembre de 2025.
        \item[Duración] 3 horas.
    \end{description}
    \newpage


    % ------------------------------------

    \begin{ejercicio}[2 puntos]
        Demostrar que el sistema de ecuaciones
        \begin{equation*}
            \left\{\begin{array}{rcl}
                    x^2-\cos(x)^3 + \cos(x^3) & = & \pi \\
                    \nicefrac{1}{y}+e^x + e^{y^4+1}-\cos(x^4-1)-\cos(\sen(x)) &=&-\pi
            \end{array}\right.
        \end{equation*}
        tiene al menos una solución en $\mathbb{R}^2$.
    \end{ejercicio}

    \begin{ejercicio}[2 puntos]
        Determina el grupo fundamental del subespacio topológico de $\mathbb{R}^3$ que se muestra en el siguiente dibujo.
    \begin{figure}[H]
            \centering
        \begin{tikzpicture}[scale=1]
            % Segmento vertical
            \draw[thick] (0,-2) -- (0,2);
            \fill[thick] (0,-2) circle(2pt) node[] {};
            \fill[thick] (0,2) circle(2pt) node[] {};

            \def\a{3.5}   % semiancho (horizontal)
            \def\c{2}   % semiancho (horizontal) chico
            \def\b{2}   % semialto (vertical mitad) -> altura total 2*b = 2

            % Dibujar solo la mitad derecha
            \draw[thick] (0,-\b) .. controls (\a,-\b) and (\a,\b) .. (0,\b);
            \draw[thick] (0,-\b) .. controls (\c,-\b) and (\c,\b) .. (0,\b);

            % Dibujar solo la mitad izquierda
            \draw[thick] (0,-\b) .. controls (-\a,-\b) and (-\a,\b) .. (0,\b);
            \draw[thick] (0,-\b) .. controls (-\c,-\b) and (-\c,\b) .. (0,\b);

            \filldraw[fill=blue!30, draw=black, opacity=0.5, thick] (2.05,0) ellipse (0.55 cm and 0.15 cm);
            \filldraw[fill=blue!30, draw=black, opacity=0.5, thick] (-2.05,0) ellipse (0.55 cm and 0.15 cm);
            \draw[thick, dotted] (-1.8,1) ellipse (0.45 cm and 0.15 cm);
            \draw[thick, dotted] (1.8,1) ellipse (0.45 cm and 0.15 cm);
            \draw[thick, dotted] (1.8,-1) ellipse (0.45 cm and 0.15 cm);
            \draw[thick, dotted] (-1.8,-1) ellipse (0.45 cm and 0.15 cm);
        \end{tikzpicture}
        \end{figure}
    \end{ejercicio}

    \begin{ejercicio}[1.5 puntos]
        Sea $X$ el subespacio topológico de $\mathbb{R}^2$: $X=\mathbb{S}^1\cup \{(x,x^2)\in \mathbb{R}^2 : x\neq 0\}$. Determinar, salvo homeomorfismo, el recubridor universal de $X$.
    \end{ejercicio}

    \begin{ejercicio}[1.5 puntos]
        Sea $\pi:\tilde{X}\to X$ una aplicación recubridora entre espacios topológicos conexos y localmente arcoconexos. Sean $f_1,f_2:Y\to \tilde{X}$ dos aplicaciones continuas, siendo $Y$ un espacio conexo y localmente arcoconexo. Supogamos
        \begin{enumerate}[label=\alph*)]
            \item $\pi\circ f_1 = \pi\circ f_2$.
            \item existe $y_0\in Y$ tal que $f_1(y_0) = f_2(y_0)$.
            \item el conjunto $\{y\in Y:f_1(y)=f_2(y)\}$ es cerrado.
        \end{enumerate}
        Demostrar que $f_1=f_2$ globalmente.
    \end{ejercicio}

    \begin{ejercicio}[1.5 puntos]
        Consideremos la superficie topológica $S$ con presentación poligonal
        \begin{equation*}
            \cc{P} = \langle a,b,c,d,e,f,g,h,i,j:ae^{-1}g^{-1}ja^{-1}bg,ch^{-1}dfh,ecdifi^{-1}j^{-1}b^{-1} \rangle .
        \end{equation*}
        Demostrar que $S$ es conexa y determinar a qué superficie modelo es homeomorfa.
    \end{ejercicio}
        
    \begin{ejercicio}
        Demostrar usando transformaciones elementales que la suma conexa $\bb{T}\# \mathbb{T} \# \mathbb{R}\mathbb{P}^2$ es homeomorfa a $\mathbb{T}\# \mathbb{R}\bb{P}^2\# \mathbb{R}\bb{P}^2\# \mathbb{R}\bb{P}^2$, siendo $\mathbb{T}$ un toro y $\mathbb{R}\bb{P}^2$ un plano proyectivo.
    \end{ejercicio}

    \newpage
    \setcounter{ejercicio}{0} % Reiniciar contador de ejercicios
    \noindent
    \textbf{Solución.}

    \begin{ejercicio}[2 puntos]
        Demostrar que el sistema de ecuaciones
        \begin{equation*}
            \left\{\begin{array}{rcl}
                    x^2-\cos(x)^3 + \cos(x^3) & = & \pi \\
                    \nicefrac{1}{y}+e^x + e^{y^4+1}-\cos(x^4-1)-\cos(\sen(x)) &=&-\pi
            \end{array}\right.
        \end{equation*}
        tiene al menos una solución en $\mathbb{R}^2$.\\

        El sistema anterior es equivalente al siguiente:
        \begin{equation*}
            \left\{\begin{array}{rcl}
                    x^2&=&\cos(x)^3 - \cos(x^3) + \pi \\
                    \nicefrac{1}{y}&=& -e^x - e^{y^4+1}+\cos(x^4-1)+\cos(\sen(x)) -\pi
            \end{array}\right.
        \end{equation*}

        Además, cualquier solución del siguiente sistema de ecuaciones también es solución del sistema anterior. Notemos que el recíproco no es necesariamente cierto, pues las soluciones con valor negativo de $x$ dejarán de serlo al tomar la raíz positiva. Sin embargo, esto no es un problema, pues si el sistema de ecuaciones que vamos a estudiar tiene solución, entonces el sistema original también la tendrá.
        \begin{equation*}
            \left\{\begin{array}{rcl}
                    x&=&\sqrt{\cos(x)^3 - \cos(x^3) + \pi} \\
                    y&=& \dfrac{1}{-e^x - e^{y^4+1}+\cos(x^4-1)+\cos(\sen(x)) -\pi}
            \end{array}\right.
        \end{equation*}

        Además, como este sistema de ecuaciones no puede tener solución con $y=0$, seguirá siendo válida en el sistema de ecuaciones original. Por lo tanto, si encontramos una solución del sistema anterior, habremos encontrado una solución del sistema original. Definimos la función $F:\mathbb{R}^2\to \mathbb{R}^2$ por
        \begin{equation*}
            F(x,y) = \left(\sqrt{\cos(x)^3 - \cos(x^3) + \pi}, \dfrac{1}{-e^x - e^{y^4+1}+\cos(x^4-1)+\cos(\sen(x)) -\pi}\right).
        \end{equation*}

        Con el objetivo de aplicar el teorema del punto fijo de Brouwer, en primer lugar hemos de comprobar que dicha función está bien definida.
        \begin{itemize}
            \item Para comprobar que $F_1$ está bien definida, basta con ver que el radicando es no negativo. Esto es cierto, pues:
            \begin{equation*}
                \cos(x)^3 - \cos(x^3) + \pi \geq -1 - 1 + \pi = \pi - 2 > 0.
            \end{equation*}
            \item Para comprobar que $F_2$ está bien definida, basta con ver que el denominador es no nulo. Esto es cierto, pues:
            \begin{equation*}
                -e^x - e^{y^4+1}+\cos(x^4-1)+\cos(\sen(x)) -\pi \leq 0-e+1+1-\pi = 2-\pi-e < -3<0
            \end{equation*}
        \end{itemize}

        Por lo tanto, $F$ está bien definida. Además, $F$ es continua, pues es la composición de funciones continuas. Ahora, vamos a acotar cada una de las componentes de $F$ para poder aplicar el teorema del punto fijo de Brouwer.
        \begin{itemize}
            \item Para la primera componente, notemos que
            \begin{equation*}
                \cos(x)^3 - \cos(x^3) + \pi \leq 1 - (-1) + \pi = 2 + \pi.
            \end{equation*}
            Por lo tanto, y por ser la raíz cuadrada una función creciente, la primera componente de $F$ está acotada superiormente por $\sqrt{2+\pi}$.
            \item Para la segunda componente, empleando que el denominador es menor que $-3$, obtenemos que:
            \begin{equation*}
                |F_2(x,y)| = \left|\dfrac{1}{-e^x - e^{y^4+1}+\cos(x^4-1)+\cos(\sen(x)) -\pi}\right| < \dfrac{1}{3}.
            \end{equation*}
        \end{itemize}

        Por tanto, como ambas componentes de $F$ están acotadas, $F$ es acotada, luego existe $R\in \mathbb{R}^+$ tal que $F(\mathbb{R}^2)\subseteq \ol{B}(0,R)$. Restringiendo $F$ a $\ol{B}(0,R)$, obtenemos una función continua $F\Big|_{\ol{B}(0,R)}:\ol{B}(0,R)\to \ol{B}(0,R)$, y por el teorema del punto fijo de Brouwer, existe $(x_0,y_0)\in \ol{B}(0,R)$ tal que $F(x_0,y_0)=(x_0,y_0)$. Por lo tanto, $(x_0,y_0)$ es solución del sistema de ecuaciones anterior, y por lo tanto también del sistema de ecuaciones original.
    \end{ejercicio}

    \begin{ejercicio}[2 puntos]
        Determina el grupo fundamental del subespacio topológico $X$ de $\mathbb{R}^3$ que se muestra en la Figura~\ref{fig:X}.
    \begin{figure}[h]
            \centering
        \begin{tikzpicture}[scale=1]
            % Segmento vertical
            \draw[thick] (0,-2) -- (0,2);
            \fill[thick] (0,-2) circle(2pt) node[] {};
            \fill[thick] (0,2) circle(2pt) node[] {};

            \def\a{3.5}   % semiancho (horizontal)
            \def\c{2}   % semiancho (horizontal) chico
            \def\b{2}   % semialto (vertical mitad) -> altura total 2*b = 2

            % Dibujar solo la mitad derecha
            \draw[thick] (0,-\b) .. controls (\a,-\b) and (\a,\b) .. (0,\b);
            \draw[thick] (0,-\b) .. controls (\c,-\b) and (\c,\b) .. (0,\b);

            % Dibujar solo la mitad izquierda
            \draw[thick] (0,-\b) .. controls (-\a,-\b) and (-\a,\b) .. (0,\b);
            \draw[thick] (0,-\b) .. controls (-\c,-\b) and (-\c,\b) .. (0,\b);

            \filldraw[fill=blue!30, draw=black, opacity=0.5, thick] (2.05,0) ellipse (0.55 cm and 0.15 cm);
            \filldraw[fill=blue!30, draw=black, opacity=0.5, thick] (-2.05,0) ellipse (0.55 cm and 0.15 cm);
            \draw[thick, dotted] (-1.8,1) ellipse (0.45 cm and 0.15 cm);
            \draw[thick, dotted] (1.8,1) ellipse (0.45 cm and 0.15 cm);
            \draw[thick, dotted] (1.8,-1) ellipse (0.45 cm and 0.15 cm);
            \draw[thick, dotted] (-1.8,-1) ellipse (0.45 cm and 0.15 cm);
        \end{tikzpicture}
        \caption{Subespacio topológico $X$ de $\mathbb{R}^3$.}
        \label{fig:X}
        \end{figure}

        En vistas de tener una notación más cómoda, sea $L$ el segmento vertical, con $N,S$ los puntos norte y sur del segmento vertical, $D_1,D_2$ los discos azules cerrados de la izquierda y derecha, respectivamente, y $H_1$, $H_2$ los ``tubos'' de la izquierda y derecha, respectivamente (que incluyen la frontera de los discos azules). Consideramos los siguientes conjuntos:
        \begin{equation*}
            U=X\setminus D_1, \quad V=X\setminus D_2.
        \end{equation*}

        Es fácil ver que $U$ y $V$ son abiertos en $X$, y que $U\cup V = X$.
        \begin{itemize}
            \item Estudiemos $U$. Al haberle eliminado el disco azul $D_1$ de la izquierda, podemos retraer la mitad superior de $H_1$ a $N$ y la mitad inferior a $S$, por lo que el conjunto $L\cup H_2\cup D_2$ es un retracto de deformación de $U$. Este conjunto es homeomorfo al representado en la Figura~\ref{fig:retracto}. Notémoslo por $R$.
            \begin{figure}[h]
                \centering
                \begin{tikzpicture}[scale=1]
                    \def\R{2} % Radio de la esfera (coincide con \b=2)

                    % 1. Sombreado sutil de la esfera (opcional, da sensación de volumen 3D)
                    \shade[ball color=gray!10, opacity=0.3] (0,0) circle (\R);

                    % 2. Ecuador relleno en azul (mismo estilo que tus discos D_1 y D_2)
                    \fill[blue!30, opacity=0.5] (0,0) ellipse (\R cm and 0.45cm);
                    % Parte trasera del ecuador (punteada)
                    \draw[thick, dotted] (\R,0) arc (0:180:\R cm and 0.45cm);
                    % Parte delantera del ecuador (continua)
                    \draw[thick] (-\R,0) arc (180:360:\R cm and 0.45cm);

                    % 3. Contorno exterior de la esfera
                    \draw[thick] (0,0) circle (\R);

                    % 4. Lazo exterior del polo norte al polo sur (corresponde al segmento L)
                    \draw[thick] (0,\R) .. controls (-4,\R) and (-4,-\R) .. (0,-\R);

                    % 5. Polos Norte y Sur
                    \fill (0,\R) circle (2pt) node[above] {$N$};
                    \fill (0,-\R) circle (2pt) node[below] {$S$};


                    % Puntos p y q
                    \fill (0,0) circle (2pt) node[left] {$p$};
                    \fill (-3,0) circle (2pt) node[left] {$q$};
                \end{tikzpicture}
                \caption{Conjunto $R$, homeomorfo a $L\cup H_2\cup D_2$.}
                \label{fig:retracto}
            \end{figure}

            Para estudiar el conjunto $R$, buscamos aplicar de nuevo el teorema de Seifert-Van Kampen. Sean:
            \begin{equation*}
                U_1 = R\setminus \{p\}, \quad V_1 = R\setminus \{q\}
            \end{equation*}
            donde $p$ pertenece al interior del disco azul y $q$ pertenece al interior del lazo exterior. Tenemos que $R=U_1\cup V_1$, con $U_1$ y $V_1$ abiertos en $R$. \begin{itemize}
                \item $U_1$ tiene como retracto de deformación la esfera junto con el lazo exterior, que sabemos que es arcoconexo con grupo fundamental isomorfo a $\bb{Z}$.
                \item $V_1$ tiene como retracto de deformación la esfera junto con el disco azul, que sabemos que es arcoconexo con grupo fundamental trivial.
                \item $U_1\cap V_1$ tiene como retracto de deformación la esfera, que sabemos que es arcoconexa con grupo fundamental trivial.
            \end{itemize}

            Por el teorema de Seifert-Van Kampen, obtenemos que el grupo fundamental de $R$ es isomorfo a $\bb{Z}$. Además, como se trata de la unión de dos conjuntos arcoconexos con intersección no vacía, $R$ es arcoconexo. Un posible generador del grupo fundamental de $R$ es la clase del lazo $\alpha$ que recorre el lazo exterior una vez, como se muestra en la Figura~\ref{fig:gen_retracto}.
            \begin{figure}[h]
                \centering
                \begin{tikzpicture}[scale=1]
                    \def\R{2} % Radio de la esfera (coincide con \b=2)

                    % 1. Sombreado sutil de la esfera (opcional, da sensación de volumen 3D)
                    \shade[ball color=gray!10, opacity=0.3] (0,0) circle (\R);

                    % 2. Ecuador relleno en azul (mismo estilo que tus discos D_1 y D_2)
                    \fill[blue!30, opacity=0.5] (0,0) ellipse (\R cm and 0.45cm);
                    % Parte trasera del ecuador (punteada)
                    \draw[thick, dotted] (\R,0) arc (0:180:\R cm and 0.45cm);
                    % Parte delantera del ecuador (continua)
                    \draw[thick] (-\R,0) arc (180:360:\R cm and 0.45cm);

                    % 3. Contorno exterior de la esfera
                    \draw[thick] (0,0) circle (\R);



                    

                    % 4. Lazo exterior del polo norte al polo sur (corresponde al segmento L)
                    \draw[thick] (0,\R) .. controls (-4,\R) and (-4,-\R) .. (0,-\R);

                    % 5. Polos Norte y Sur
                    \fill (0,\R) circle (2pt) node[above] {$N$};
                    \fill (0,-\R) circle (2pt) node[below] {$S$};

                    % --- GENERADOR DEL GRUPO FUNDAMENTAL (en rojo) ---
                    % Bajada por el lazo exterior (de N a S)
                    \draw[very thick, red, postaction={decorate}, decoration={markings, mark=at position 0.5 with {\arrow{>}}}, ->]
                        (0,\R) .. controls (-4,\R) and (-4,-\R) .. (0,-\R);

                    % Subida por la superficie de la esfera (de S a N)
                    \draw[very thick, red, postaction={decorate},
      decoration={markings, mark=at position 0.5 with {\arrow{>}}}]
      (0,-\R) arc (-90:90:\R);


                    % Puntos p y q
                    \fill (0,0) circle (2pt) node[left] {$p$};
                    \fill (-3,0) circle (2pt) node[left] {$q$};

                    % Etiqueta del lazo generador
                    \node at (-2.5,1.5) {\textcolor{red}{$\alpha$}};

                \end{tikzpicture}
                \caption{Generador del grupo fundamental de $R$.}
                \label{fig:gen_retracto}
            \end{figure}

            Por lo tanto, $U$ es arcoconexo y su grupo fundamental es isomorfo a $\bb{Z}$.


            \item Estudiemos $V$. Como $V\cong U$ (por simetría), obtenemos que $V$ es arcoconexo y su grupo fundamental es isomorfo a $\bb{Z}$.


            \item Estudiemos $U\cap V$. Al haber eliminado los discos azules, podemos retraer la mitad superior de cada $H_i$ a $N$ y la mitad inferior a $S$, por lo que $L$ es un retracto de deformación de $U\cap V$. Por lo tanto, $U\cap V$ es arcoconexo y su grupo fundamental es trivial.
        \end{itemize}

        Por el teorema de Seifert-Van Kampen, obtenemos que el grupo fundamental de $X$ es isomorfo a $\bb{Z}\ast \bb{Z}$, y además $X$ es arcoconexo. Su generador puede ser la clase del lazo $\alpha'$ que recorre el lazo exterior de la derecha una vez y la clase del lazo $\beta'$ que recorre el lazo exterior de la izquierda una vez, como se muestra en la Figura~\ref{fig:gen_inicial}.
        \begin{figure}[h]
            \centering
            \begin{tikzpicture}[scale=1]
                \def\a{3.5}   % semiancho (horizontal)
                \def\c{2}     % semiancho (horizontal) chico
                \def\b{2}     % semialto (vertical mitad)

                % Contorno mitad derecha (H_2)
                \draw[thick] (0,-\b) .. controls (\a,-\b) and (\a,\b) .. (0,\b);
                \draw[thick] (0,-\b) .. controls (\c,-\b) and (\c,\b) .. (0,\b);

                % Contorno mitad izquierda (H_1)
                \draw[thick] (0,-\b) .. controls (-\a,-\b) and (-\a,\b) .. (0,\b);
                \draw[thick] (0,-\b) .. controls (-\c,-\b) and (-\c,\b) .. (0,\b);

                % Discos azules D_1 y D_2
                \filldraw[fill=blue!30, draw=black, opacity=0.5, thick] (2.05,0) ellipse (0.55 cm and 0.15 cm);
                \filldraw[fill=blue!30, draw=black, opacity=0.5, thick] (-2.05,0) ellipse (0.55 cm and 0.15 cm);

                % Secciones punteadas
                \draw[thick, dotted] (-1.8,1) ellipse (0.45 cm and 0.15 cm);
                \draw[thick, dotted] (1.8,1) ellipse (0.45 cm and 0.15 cm);
                \draw[thick, dotted] (1.8,-1) ellipse (0.45 cm and 0.15 cm);
                \draw[thick, dotted] (-1.8,-1) ellipse (0.45 cm and 0.15 cm);

                % --- GENERADOR DEL GRUPO FUNDAMENTAL (en rojo) ---
                % Bajada por el segmento vertical L (de N a S)
                \draw[very thick, red, postaction={decorate},
                    decoration={markings, mark=at position 0.3 with {\arrow{>}}}]
                    (0.02,\b) -- (0.02,-\b);

                % Subida por la pared exterior derecha de H_2 (de S a N)
                \draw[very thick, red, postaction={decorate},
                    decoration={markings, mark=at position 0.5 with {\arrow{>}}}]
                    (0,-\b) .. controls (\a,-\b) and (\a,\b) .. (0,\b);


                % Bajada por el segmento vertical L (de N a S)
                \draw[very thick, green!30!black, postaction={decorate},
                    decoration={markings, mark=at position 0.8 with {\arrow{>}}}]
                    (-0.02,\b) -- (-0.02,-\b);
                
                    % Subida por la pared exterior izquierda de H_1 (de S a N)
                \draw[very thick, green!30!black, postaction={decorate},
                    decoration={markings, mark=at position 0.5 with {\arrow{>}}}]
                    (0,-\b) .. controls (-\a,-\b) and (-\a,\b) .. (0,\b);

                % Polos Norte y Sur
                \fill (0,\b) circle (2pt) node[above] {$N$};
                \fill (0,-\b) circle (2pt) node[below] {$S$};

                % Etiqueta del lazo generador
                \node at (2.8,1.3) {\textcolor{red}{$\alpha$}};

                \node at (-2.8,1.3) {\textcolor{green!30!black}{$\beta$}};

            \end{tikzpicture}
            \caption{Generadores del grupo fundamental de $X$.}
            \label{fig:gen_inicial}
        \end{figure}
    \end{ejercicio}

    \begin{ejercicio}[1.5 puntos]
        Sea $X$ el subespacio topológico de $\mathbb{R}^2$: $X=\mathbb{S}^1\cup \{(x,x^2)\in \mathbb{R}^2 : x\neq 0\}$. Determinar, salvo homeomorfismo, el recubridor universal de $X$.
    \end{ejercicio}

    \begin{ejercicio}[1.5 puntos]
        Sea $\pi:\tilde{X}\to X$ una aplicación recubridora entre espacios topológicos conexos y localmente arcoconexos. Sean $f_1,f_2:Y\to \tilde{X}$ dos aplicaciones continuas, siendo $Y$ un espacio conexo y localmente arcoconexo. Supogamos
        \begin{enumerate}[label=\alph*)]
            \item $\pi\circ f_1 = \pi\circ f_2$.
            \item existe $y_0\in Y$ tal que $f_1(y_0) = f_2(y_0)$.
            \item el conjunto $\{y\in Y:f_1(y)=f_2(y)\}$ es cerrado.
        \end{enumerate}
        Demostrar que $f_1=f_2$ globalmente.
    \end{ejercicio}

    \begin{ejercicio}[1.5 puntos]
        Consideremos la superficie topológica $S$ con presentación poligonal
        \begin{equation*}
            \cc{P} = \langle a,b,c,d,e,f,g,h,i,j:ae^{-1}g^{-1}ja^{-1}bg,ch^{-1}dfh,ecdifi^{-1}j^{-1}b^{-1} \rangle .
        \end{equation*}
        Demostrar que $S$ es conexa y determinar a qué superficie modelo es homeomorfa.
    \end{ejercicio}
        
    \begin{ejercicio}
        Demostrar usando transformaciones elementales que la suma conexa $\bb{T}\# \mathbb{T} \# \mathbb{R}\mathbb{P}^2$ es homeomorfa a $\mathbb{T}\# \mathbb{R}\bb{P}^2\# \mathbb{R}\bb{P}^2\# \mathbb{R}\bb{P}^2$, siendo $\mathbb{T}$ un toro y $\mathbb{R}\bb{P}^2$ un plano proyectivo.
    \end{ejercicio}


\end{document}
